\section*{第 \S\arabic{exercisegroup} 节习题}
\phantomsection
\addcontentsline{toc}{section}{第 \protect\S\arabic{exercisegroup} 节习题}

\begin{enumerate}
\def\labelenumi{\arabic{enumi})}
\tightlist
\item
  设 \(\alpha\) 与 \(\beta\) 是两个有限实数，满足 \(\alpha < \beta\)。证明：若 \(\gamma\) 与 \(\delta\) 是两个数，满足 \(\alpha < \gamma \leqslant \delta < \beta\)，则存在实函数 \(f\)（定义在区间 \([0, a]\) 上），只取值 \(\alpha\) 与 \(\beta\)，使得在整个区间 \([\varepsilon, a]\)（对 \(\varepsilon > 0\)）上 \(f\) 是阶梯函数，并且若令 \(g(x) = \int_0^x f(t)dt\)，则
\end{enumerate}

\[
\liminf _ {x \to 0,   x > 0} \frac {g (x)}{x} = \gamma , \quad \text { and } \quad \limsup _ {x \to 0,   x > 0} \frac {g (x)}{x} = \delta ;\tag{*}
\]

（取 g 为这样的函数：其图象是一条折线，相继各边的斜率为 \(\alpha\) 与 \(\beta\)，其峰点交替地落在直线 \(y = \gamma x\) 与 \(y = \delta x\) 上（当 \(\gamma \neq \delta\)），或落在直线 \(y = \gamma x\) 与抛物线 \(y = \gamma x + x^{2}\) 上（当 \(\gamma = \delta\)））。

用同样的方法证明：不论 \(\gamma\) 与 \(\delta\) 有限与否（\(\gamma < \delta\)），都存在定义在 \([0, a]\) 上的实函数 f，使得在每个区间 \([\varepsilon, a]\)（对 \(\varepsilon > 0\)）上 f 是阶梯函数，积分 \(g(x) = \int_{0}^{x} f(t) dt\) 存在，并且仍然有关系式 (*)。

\begin{enumerate}
\def\labelenumi{\arabic{enumi})}
\setcounter{enumi}{1}
\tightlist
\item
  \begin{enumerate}
  \def\labelenumii{\alph{enumii})}
  \tightlist
  \item
    设 \(\mathbf{f}\) 是区间 \([0, a]\) 上的规则函数，使得积分 \(\int_0^a \frac{\mathbf{f}(t)}{t} dt\) 收敛。证明积分 \(\mathbf{g}(x) = \int_0^x \mathbf{f}(t) dt\) 收敛，且 \(\mathbf{g}\) 在点 \(x = 0\) 处有右导数（适当地分部积分）。
  \end{enumerate}
\end{enumerate}

\begin{enumerate}
\def\labelenumi{\alph{enumi})}
\setcounter{enumi}{1}
\tightlist
\item
  给出一个实函数 \(f\) 的例子，使得积分 \(g(x) = \int_0^x f(t)dt\) 收敛并且在点 0 处有右导数，然而 \(f(t) / t\) 在 \([0, a]\) 上没有积分（取 \(f(x)/x\) 为形如 \(\cos\varphi(x)\) 的函数的导数，其中 \(\varphi\) 趋于 \(+\infty\)（当 x 趋于 0 时））。
\end{enumerate}

\begin{enumerate}
\def\labelenumi{\arabic{enumi})}
\setcounter{enumi}{2}
\item
  设 \(f\) 是实函数，\(\geqslant 0\)，定义在区间 \(]0, a]\) 上且在这个区间上是规则的，使得积分 \(g(x) = \int_0^x f(t)dt\) 收敛，但 \(g\) 在点 0 处没有右导数。证明：存在一个函数 \(f_1\)，在 \(]0, a]\) 上是规则的，使得 \((f_1(x))^2 = f(x)\)（对一切 \(x\)），积分 \(g_1(x) = \int_0^x f_1(t)dt\) 收敛，且 \(g_1\) 在 0 处有右导数。（首先考虑 \(f\) 在每个区间 \([\varepsilon, a]\)（对 \(\varepsilon > 0\)）上等同于一个阶梯函数的情形。在这种情形，把 \(f\) 取常值的每个区间分成足够多等份，并取 \(f_1\) 在这些区间的每一个上取常值，\(f_1\) 的符号在两个相邻的这种区间上不同。一般情形同样处理。）
\item
  在区间 \([0,1]\) 上定义两个实函数 f, g，使得 f 与 g 有严格原函数，但 fg 在 \([0,1[\) 的每一点处都是一个连续函数的右导数，而这个连续函数在一个具有连续统的势的点集上没有左导数（利用与 I, p.~38 习题 8 及前面习题 3 类似的构造）。
\item
  \begin{enumerate}
  \def\labelenumii{\alph{enumii})}
  \tightlist
  \item
    设 \(\mathbf{f}\) 是有界开区间 \(]a\) , \(b[\) 上的规则函数；假设存在实函数 \(g\)，在 \(]a\) , \(b[\) 上递减，使得 \(\| \mathbf{f}(x) \| \leqslant g(x)\)（在 \(]a\) , \(b[\) 上），且积分 \(\int_{a}^{b} g(t) dt\) 收敛。证明：若 \((\varepsilon_n)\) 是一个数列，\(>0\)、趋于 0，且满足 \(\inf_{n > 1} n \varepsilon_n > 0\)，则
  \end{enumerate}
\end{enumerate}

\[
\lim _ {n \rightarrow + \infty} \frac {1}{n} \sum_ {k = 1} ^ {n - 1} \mathbf {f} \left(a + \varepsilon_ {n} + k \frac {b - a}{n}\right) = \int_ {a} ^ {b} \mathbf {f} (t) d t.\tag{*}
\]

\begin{enumerate}
\def\labelenumi{\alph{enumi})}
\setcounter{enumi}{1}
\item
  给出一个实规则函数 f（\textgreater{} 0，在 \([0,1]\) 上）的例子，使积分 \(\int_{0}^{1} f(t) dt\) 收敛，而关系式 (*) 对 \(\varepsilon_{n} = 1/n\) 不成立（取 f 使它在 \(x = 2^{-p}\) 处的值为 \(2^{2p}\)）。
\item
  在与 a) 相同的假设下，证明
\end{enumerate}

\[
\lim _ {n \rightarrow + \infty} \frac {1}{n} \sum_ {k = 1} ^ {n} (- 1) ^ {k} \mathbf {f} \left(a + k \frac {b - a}{n}\right) = 0.
\]

\begin{enumerate}
\def\labelenumi{\arabic{enumi})}
\setcounter{enumi}{5}
\tightlist
\item
  设 \(\mathbf{f}\) 是区间 \(]a, +\infty[\) 上的规则函数；假设存在实递减函数 \(g\)（在 \(]a, +\infty[\) 上），使得 \(\| \mathbf{f}(x)\| \leqslant g(x)\)（在这个区间上），且积分 \(\int_{a}^{+\infty}g(t)dt\) 收敛。证明：级数 \(\sum_{n=1}^{\infty}\mathbf{f}(a + nh)\) 对每个 \(h > 0\) 绝对收敛，且
\end{enumerate}

\[
\lim _ {h \rightarrow + \infty} h \sum_ {n = 1} ^ {\infty} \mathbf {f} (a + n h) = \int_ {a} ^ {+ \infty} \mathbf {f} (t) d t.
\]

\begin{enumerate}
\def\labelenumi{\arabic{enumi})}
\setcounter{enumi}{6}
\item
  设 \(f\) 与 \(g\) 是两个规则函数，且 \(>0\)（在开区间 \(]a, b[\) 上）。证明：函数 \(f / (1 + fg)\) 与 \(\inf(f, 1/g)\) 在 \(]a, b[\) 上的积分同时收敛或同时为无穷。
\item
  设 \(f\) 是规则函数且 \(\geqslant 0\)（在区间 \([a, +\infty[\) 上），\(g\) 是定义在 \([a, +\infty[\) 上的可微递增函数，且满足 \(g(x) - x \geqslant \lambda > 0\)（对所有 \(x \geqslant a\)）。
\end{enumerate}

证明：若有 \(f(g(x))g'(x) \leqslant kf(x)\)（其中 k \textless{} 1）（相应地，\(f(g(x))g'(x) \geqslant kf(x)\)（其中 k \textgreater{} 1）），则积分 \(\int_{a}^{+\infty} f(t) dt\) 收敛（相应地，等于 \(+\infty\)）（叶尔马科夫判别法：把 g 的 \(n^{th}\) 次迭代记作 \(g^n\)，考察积分 \(\int_{0}^{g^n(a)} f(t) dt\) 并让 n 无限增大）。

\begin{enumerate}
\def\labelenumi{\arabic{enumi})}
\setcounter{enumi}{8}
\tightlist
\item
  设 \(\alpha\) 是一个数 \(>0\)，\(f\) 是定义在区间 \([0, +\infty[, \geqslant 0\) 且递减的函数，使得积分 \(\int_0^{+\infty} t^\alpha f(t) dt\) 收敛。证明：对所有 \(x > 0\)，有
\end{enumerate}

\[
\int_ {x} ^ {+ \infty} f (t) d t \leqslant \left(\frac {\alpha}{(\alpha + 1) x}\right) ^ {\alpha} \int_ {0} ^ {+ \infty} t ^ {\alpha} f (t) d t.
\]

（先对 \(f\) 在区间 \(]0, a]\) 上取常值且对 \(x > a\) 为零的情形证明，再对这种函数之和证明，一般情形通过取极限得到。）
% semantic-fragments: legacy-input-seam
\relax{}
\setcounter{exercisegroup}{2}
\refstepcounter{exercisegroup}
